% ===========================================================================
% VI. Discussion
% ===========================================================================
\section{Discussion}
\label{sec:discussion}

\subsection{Defense Mechanism Effectiveness}

The three defense classes exhibit qualitatively different effectiveness
profiles under MARL adversarial training.

\textbf{Deterministic defenses} (self-rating, admin escalation, gate bypass)
achieve 100\% block rates because the negative expected return is
unconditional: no matter how many times an adversarial agent attempts these
actions, the penalty always exceeds the reward bonus.  Under MARL training,
this means that the gradient signal is uniformly negative for these actions,
and the policy converges to never selecting them.  The one caveat is that
the policy may still assign non-zero probability to these actions at the
local equilibrium (observed in C6's $93.2\%$ honest rate), reflecting that
the entropy coefficient $c_e = 0.05$ keeps the policy stochastic even at
convergence.  Reducing $c_e$ at late training stages could eliminate this
residual.

\textbf{Probabilistic detection} (evidence tampering, $p=0.80$) produces
a nuanced equilibrium.  The negative expected return ($\mathbb{E}[R] = -1.4$
per attempt) drives convergence toward honest behavior, but the nonzero
evasion probability ($20\%$) creates a reward signal variance that prevents
full convergence.  In a real deployment, the detection probability could be
increased (e.g., by requiring on-chain cryptographic commitments for all
evidence) to push the equilibrium further toward deterrence.

\textbf{Economic deterrence} (Sybil, collusion) provides only $52$--$58\%$
block rates.  This is by design: economic mechanisms do not prevent attacks
but impose costs that limit their profitability.  The honest population
still converges to $\geq\!87\%$ honest behavior even under Sybil pressure
(C3: $87.6\%$), confirming that economic deterrence reduces attack frequency at the
population level, though the observed $52$--$58\%$ block rates indicate that
these mechanisms remain \emph{persistently exploitable} under adaptive
adversarial strategies.  Economic deterrence imposes costs that shift the
equilibrium toward honesty but does not eliminate exploitation channels:
approximately half of all Sybil and collusion attempts succeed, and a
sufficiently capitalized adversary could sustain this exploitation
indefinitely.  Tightening the stake cost or CI widening rate would shift
this equilibrium further.

\subsection{Noop Equilibrium and Rational Abstention}

An important observation is that all agents — including adversarial agents
with reward bonuses up to $2\times$ — ultimately converge to \emph{rational
abstention}: preferring no-operation (A0) over risky attack actions when
the expected return is negative.  This is consistent with the game-theoretic
prediction that rational agents will not repeatedly attempt actions with
negative expected value, regardless of the reward bonus magnitude.  The
MARL training validates that this theoretical prediction holds empirically
for all single-attack configurations.

We note an important distinction between \emph{security robustness} and
\emph{functional utility}.  A population that converges to no-operation
(A0) rather than honest rating (A1--A2) is safe from a security perspective ---
no attacks succeed --- but does not achieve the system's intended purpose of
producing accurate, evidence-grounded reputation scores.  The honest action
percentage metric captures this distinction: configurations where agents
converge to honest \emph{rating} (C1, C4, C7--C9: $\geq\!99.5\%$ honest)
provide both security and utility, while configurations where agents converge
to abstention provide security without utility.  Designing reward structures
that incentivize active honest participation --- not merely attack avoidance ---
is an open challenge that the terminal reward result (Section~\ref{sec:terminal})
begins to address.

\subsection{Multi-Vector Attack as the Open Challenge}

The C11 result ($20.0\% \pm 40.0\%$ honest, 4/5 seeds adversarial-dominant)
indicates that the honest Nash equilibrium, while stable once reached, is
\emph{not globally attracting} at 50\% adversarial fraction.  The system
possesses two stable attractors --- the honest equilibrium and the adversarial
equilibrium --- with the basin of attraction depending on early training
dynamics.  In a real deployment, this translates to \emph{path dependence}:
an early-stage coordinated attack during system bootstrap could lock the
network into the adversarial basin before honest participants accumulate
sufficient evidence mass.  The terminal reward result
(Section~\ref{sec:terminal}) demonstrates that this path dependence is
addressable through reward design rather than requiring architectural changes
to the defense layer.

The comprehensive configuration (C11) reveals a qualitatively different
challenge.  With $50\%$ adversarial population and coordinated multi-vector
attacks, the interaction between attack vectors creates emergent dynamics
that no single-attack defense fully handles.  Specifically:
\begin{itemize}[noitemsep]
  \item Sybil identities can attempt gate bypass with the Sybil's reduced
        stake, creating a compound attack that neither the Sybil defense
        nor the gate bypass defense fully handles in isolation.
  \item Colluding agents can use provenance replay to amplify positive
        ratings for colluding partners, combining two partially-defended
        attack vectors.
  \item The $50\%$ adversarial fraction means that even the defense dividend
        ($+0.1$ to honest agents per blocked attack) cannot overcome the
        adversarial reward accumulation from unblocked Sybil and colluding
        actions.
\end{itemize}
Fundamentally, the C11 failure arises from the inability of honest agents
to \emph{coordinate} their responses against adversarial clusters: each
honest agent optimizes independently based on its own 14-dimensional
observation, with no mechanism to detect, signal, or collectively respond
to coordinated multi-vector attacks.  This is not merely a useful
extension but the \emph{primary structural limitation} of the current
framework.

Addressing C11 likely requires extending the observation space with
cross-agent comparison (e.g., including the distribution of stake levels
across all agents) and introducing a cooperative honest majority protocol
that enables honest agents to identify and coordinate against adversarial
clusters.

\subsection{Terminal Reward as a Design Recommendation}
\label{sec:terminal_design}

The terminal reward variant (Section~\ref{sec:terminal}) produces the most
striking result in the entire evaluation: C11 honest behavior improves from
$20.0\%$ to $100.0\%$ across all 3 seeds when a reputation-progress component
$R_{\text{terminal}} = 0.5 \times (\mu_i^{\text{final}} - 0.5)$ is added at
episode end.  This result demonstrates that the C11 failure is not a
fundamental limitation of the defense mechanism design but rather a consequence
of the reward objective: step-level action rewards do not capture the long-term
cost of reputation degradation under multi-vector coordination.
More broadly, this exposes a fundamental limitation of step-wise reward
design in adversarial multi-agent settings: local action-level incentives
are insufficient to enforce global system-level objectives when adversaries
can coordinate across multiple attack vectors simultaneously.  The step
reward treats each action in isolation, but the system's integrity is a
\emph{trajectory-level} property that emerges from the cumulative
interaction of all agents' actions over the full episode.

The asymmetry between C2 and C11 is instructive.  For C2 (5 adversaries,
$b = 0.5$), the terminal reward does not eliminate bimodality
($65.0\% \pm 46.0\%$ honest), because the lower adversarial pressure and
fewer attack vectors make the per-step reward landscape already navigable ---
the terminal signal adds insufficient gradient to escape the adversarial basin.
For C11 (10 adversaries, $b = 1.0$, all 9 attack vectors), the cumulative
reputation damage from coordinated multi-vector attacks is large enough that
the terminal penalty creates a steep gradient away from the adversarial
equilibrium, collapsing the bimodal landscape into a single honest attractor.

This finding has a direct design implication for deployed reputation systems:
the RL evaluation suggests that \emph{outcome-based incentive alignment} ---
connecting agent rewards to long-term reputation consequences, not just
per-action penalties --- is essential for robust defense under coordinated
multi-vector attacks.  In the blockchain implementation, this translates to
designing economic mechanisms where an actor's cumulative reputation
trajectory, not just individual transaction outcomes, determines their
net economic return.

\subsection{Methodology Comparison with Scripted Evaluation}

Our MARL results are broadly consistent with the scripted evaluation in the
companion paper~\cite{almaayn2024unified}: both find 100\% block rates for
self-rating, admin escalation, and gate bypass; both find partial mitigation
for Sybil and collusion; both find that evidence tampering is probabilistically
defended.  The MARL evaluation adds quantitative block rate measurements
(previously only pass/fail was reported), reveals the honest \% convergence
trajectory, and discovers the multi-vector challenge in C11.  The two methods
are therefore complementary: scripted evaluation confirms defense behavior for
specific known exploits; MARL training provides convergence guarantees and
discovers emergent multi-vector dynamics.

\subsection{Limitations}

\textbf{Simulation fidelity:} The MARL environment is a simplified model of
the blockchain system.  It abstracts away network latency, smart contract
execution overhead, MVCC conflicts, and Fabric-level permission enforcement.
Defense mechanisms are modeled as reward penalties rather than hard code
rejections; a real deployment cannot override a successful chaincode
transaction via a negative reward.

\textbf{Fixed adversary population:} The adversarial fraction is fixed at
initialization and does not change during training.  In practice, adversarial
agents may enter or leave the supply chain, and the defense system should be
evaluated under dynamic adversarial populations.

\textbf{Observation space completeness:} The 14-dimensional observation does
not include other agents' actions or reputations directly.  Including partial
observability of peer behavior could improve both honest convergence and
adversarial detection.

\subsection{Learned Avoidance vs.\ Mechanism Robustness}

An important methodological caveat applies to the interpretation of all
MARL-based robustness results.  When a trained agent learns to
\emph{avoid} an attack action due to its associated penalty, the
observed low attack rate reflects \emph{policy adaptation}, not
necessarily \emph{intrinsic defense strength}.  Provenance replay (C10)
illustrates this distinction clearly: the mechanism's block rate is only
27.1\%, yet trained agents rarely attempt it because the $-5.0$ penalty
makes every attempt unprofitable regardless of detection.  The defense
mechanism itself remains weak --- a non-RL adversary (e.g., a scripted
attacker, a human operator, or an agent trained with a different reward
structure) could still exploit the 72.9\% evasion rate.

This distinction generalizes across all configurations.  The high honest
percentages reported in Table~\ref{tab:summary} should be interpreted as
upper bounds on robustness \emph{against the specific MARL adversary class
evaluated}, not as guarantees against arbitrary adversarial strategies.
MARL evaluation complements scripted testing by discovering emergent
strategies and quantifying convergence dynamics, but it does not replace
formal verification or red-team exercises with human adversaries.  The
robustness observed in this evaluation may overestimate security against
adversaries not constrained by the learned policy class.
